% ════════════════════════════════════════════
% 清洗与编码（5.2.2）
% ════════════════════════════════════════════
\subsubsection{清洗与编码}

原始数据只有 1 条缺失、0 条重复，但类别变量的稀疏程度与因变量的重尾程度都超出线性回归的适用条件，若直接代入会导致设计矩阵病态、残差严重非正态。为使后续回归稳定可估，数据清洗按缺失处理、类别编码、稀疏时段合并与对数变换四步展开。

首先处理缺失。付费标记 X7 仅有 1 条缺失，占 0.2\%，且该帖子的其余字段完整，删除会浪费一条有效信息，故采用众数填补，取值为 0。填补只修改 X7 一列，并附加填补标记，便于追溯与敏感性对照。

随后完成编码。对帖子类型与内容主题采用独热编码，并各舍弃一个水平作为参照：帖子类型以 Photo 为参照，保留 Status、Link、Video 三个哑变量；内容主题以主题 1 为参照，保留主题 2、主题 3 两个哑变量。月份与星期分别生成 11 个和 6 个哑变量，参照类为 1 月与周一，从而避免虚拟变量陷阱。

然后合并稀疏时段。发帖小时在样本中有 22 个取值，0 时与 21 时没有帖子；若逐小时生成哑变量，部分单元只有个位数观测，估计方差极大。按样本分布将其合并为凌晨（0 至 5 时）、上午（6 至 11 时）与下午晚间（12 至 23 时）三组，并以上午为参照。合并后用两个哑变量代替原来的 21 个，牺牲少量时间分辨率，换取每个单元足够的样本量。

最后变换因变量与规模变量。曝光量为非负计数且右偏严重，采用 log1p 变换，即对曝光次数加 1 后取自然对数；加 1 保证曝光为零时有定义，本样本最小值为 570，实际不受影响。粉丝规模采用对数变换，把“每增加一人”改为“每增加 1\%”的弹性口径。两处变换都只改变尺度，不改变样本排序。

预处理的核心决策是缺失值如何处理。针对这一任务，需在行删除、均值或中位数填补、众数填补与多重插补等方案中择优，从适用变量类型、信息损失与实现成本三个维度对比其适用性，如表\ref{tab:algo2}所示。

\begin{longtable}{>{\centering\arraybackslash}p{0.18\textwidth}>{\centering\arraybackslash}p{0.20\textwidth}>{\centering\arraybackslash}p{0.24\textwidth}>{\centering\arraybackslash}p{0.26\textwidth}}
  \caption{缺失处理方法对比}
  \label{tab:algo2} \\
  \toprule
  方法 & 适用变量 & 信息损失 & 主要局限 \\
  \midrule
  \endfirsthead
  \caption{缺失处理方法对比（续）} \\
  \toprule
  方法 & 适用变量 & 信息损失 & 主要局限 \\
  \midrule
  \endhead
  \bottomrule
  \endfoot
  行删除 & 任意类型 & 损失 1 条完整记录 & 样本量减少，估计效率下降 \\
  均值或中位数 & 连续变量 & 保持中心位置 & 忽略变量间关系，压低方差 \\
  众数填补 & 离散或二元变量 & 保持峰值类别 & 会强化多数类，需做对照 \\
  多重插补 & 任意类型 & 反映填补不确定性 & 需要辅助模型，实现复杂 \\
\end{longtable}

如表\ref{tab:algo2}所示，本部分的缺失发生在二元变量 X7 上，且只缺 1 条，多重插补的收益不足以抵消其复杂度；均值填补对 0/1 变量没有意义。众数填补能保持付费标记的多数类结构，同时保留该帖子的其他字段，是当前情形下最合适的选择。为排除填补带来的偏差，本节同时估计删除缺失行的方案，若两者系数差异很小，则可认为结论不依赖填补方式。综合以上分析，本部分选择众数填补为主、删除方案为辅的策略。
